Con los resultados obtenidos se puede concluir que el comportamiento práctico de los tres algoritmos está relacionado con sus complejidades teóricas. Greedy mostró los menores tiempos de ejecución, lo que concuerda con su complejidad O(N+M), registrando aproximadamente 0,25 ms en las grillas libres de 1024×2048, mientras que Programación Dinámica, con una complejidad O(NMV), alcanzó alrededor de 95 ms. En cambio, Fuerza Bruta superó los 10 segundos en las grillas libres de 16×16. Con respecto a la memoria, Programación Dinámica presentó un mayor consumo debido a su complejidad espacial O(NMV), alcanzando aproximadamente 102.000 KiB, valor bastante alto en comparación a los 28.000 KiB de Greedy en las grillas más grandes y Fuerza Bruta presentó un consumo de memoria bajo en las instancias evaluadas, ya que no almacena todos los caminos generados. Luego, se observó que Greedy no garantiza encontrar la solución óptima e incluso puede fallar, mientras que Programación Dinámica encontró solución en las 288 instancias evaluadas y Fuerza Bruta obtuvo los mismos resultados que Programación Dinámica en las 120 instancias donde ambos algoritmos fueron ejecutados , cumpliendo así el objetivo de relacionar los resultados experimentales con las complejidades teóricas.
